% =========================
% Chapter 8
% =========================
\chapter{Limitations}
\label{ch:limitations}

We state the system's limitations plainly. Several are fundamental consequences of
the no-infrastructure design bet --- guiding a blind user with nothing but a phone
camera, a compass, and a server --- rather than implementation gaps; others are
scope decisions with clear extension paths (Chapter~\ref{ch:future_work}).

\section{Perception and Navigation}

\begin{enumerate}[nosep]
  \item \textbf{No localization.} Without SLAM, maps, or markers, the system has
        no notion of where it is in a building. Each room is explored on local
        evidence only, and the system does not remember which door it entered
        through or which doors it has already tried --- so after a doorway
        transit, the entry door will be re-found by the new room's scan, and a
        journey through a goal-less branch can, in principle, ping-pong between
        two rooms or re-enter a dead end. In practice the user's global ``stop''
        and their own coarse sense of the space bound this, but the system itself
        does not.
  \item \textbf{Direction hints are shallow.} At a multi-door fork the system
        asks, and a spoken ``left'' or ``right'' redirects the choice --- but the
        hint applies to that fork only, in the current room, and is forgotten the
        moment it is used. A hint that matches no confirmed door falls back to
        the autonomous pick (nearest straight ahead), and nothing is remembered
        across forks: the same doorway may be asked about again on a later pass,
        and a hint cannot yet be embedded in the initial command (``the kitchen
        is to the left''). Deeper hint handling and per-fork memory are staged in
        Chapter~\ref{ch:future_work}.
  \item \textbf{Domain shift can defeat the door models.} The door detector's
        strong validation metrics (Section~\ref{sec:door_eval}) measure
        generalization within its public training datasets. In some real rooms,
        textured walls fooled all three models --- the primary detector, COCO,
        and the verifier --- simultaneously. The perception funnel contains the
        damage (such candidates rarely survive the gates, and never the sighting
        minimums), but containment is not a cure; the durable fix is a
        hard-negative fine-tune on exactly such footage.
  \item \textbf{Advisory guidance, not collision avoidance.} Perception runs at
        3--5\,FPS over a network hop; its value is advance information in the
        2--4\,m band --- where the door is, what is in the path, which room this
        is. The white cane remains responsible for the immediate contact zone,
        and every trial kept a sighted spotter present. The obstacle watchdog is
        additionally armed only during the door approach --- the phase in which
        the system has instructed the user to walk --- not during scans or turns.
  \item \textbf{Monocular distance is approximate.} Door distances assume a
        2\,m door and a nominal field of view; non-standard doors, unusual
        lenses, or digital zoom shift the step counts. The tighter-box
        measurement, the median smoothing, and the tactile endgame (``reach out
        with your hand'') are designed so that the final meter never depends on
        the estimate being exact.
  \item \textbf{Compass dependence.} The $360^\circ$ scan's bearings inherit the
        magnetometer's weaknesses --- and magnetic interference is strongest
        near exactly the large appliances that define a kitchen. Compass
        glitches are filtered (large jumps ignored; completion requires the
        confirmed return to start), and a compass-less fallback exists, but the
        fallback loses bearing summaries and the guided turn.
\end{enumerate}

\section{System and Deployment}

\begin{enumerate}[nosep]
  \item \textbf{Network and cloud dependence.} All perception and speech run
        server-side: a session requires connectivity for camera streaming, and
        speech synthesis requires reachability of the online TTS service. The
        phrase cache absorbs repeated lines, and the resume mechanism recovers
        the task after a drop, but a dead connection means a silent system ---
        mitigated only by the client's local speech fallback announcing the
        disconnection.
  \item \textbf{Latency floors.} Every frame and every utterance crosses the
        network. The architecture hides this well in the steady state (worker
        threads, latest-frame reads, cached phrases), but the floor is real:
        time-critical reactions --- a person stepping into the path at close
        range --- are bounded by capture, transport, inference, synthesis, and
        playback in sequence. This reinforces, rather than merely accompanies,
        the cane-complementary scope.
  \item \textbf{English only; curated vocabularies.} Voice interaction is
        English; Object Allocation covers a curated set of COCO classes;
        Navigation recognizes six destination room types. Rooms defined by
        objects outside COCO --- hallways, staircases, garages --- cannot be
        arrival targets at all, since the arrival rule needs indicator classes
        the detector can produce.
  \item \textbf{Perception logic is validated live, not by the automated
        suite.} The pure decision layer is exhaustively tested, but the pixel
        side of the door funnel --- edge densities, IoU corroboration against
        real frames, the segmentation lane --- is validated only by the live
        trials. A regression there would be caught by a person, not by the test
        run.
  \item \textbf{Demonstration-grade hardening.} Detection debug tracing is left
        enabled by design (it is how live failures are diagnosed) and would be
        disabled for production; captured push-to-talk audio is stored on the
        server for debugging, which is a privacy consideration a deployment
        would need to address; and the deployment story is a developer tunnel,
        not a hardened multi-tenant host.
\end{enumerate}

\section{Evaluation Scope}

\begin{enumerate}[nosep]
  \item \textbf{No blind-user study.} All live trials used sighted testers.
        The system's central claims about cognitive load, trust, and usability
        are grounded in design principles from the assistive-technology
        literature and in the trials' qualitative behavior --- but validating
        them for the target population requires a formal study with blind and
        low-vision participants, under an appropriate safety protocol. This is
        the single most important item of future work.
  \item \textbf{Demonstrated scope is small-scale.} Live journeys were two rooms
        deep in residential settings with openable doors. Building-scale
        environments --- long corridors, identical doors, elevators, glass ---
        are outside the demonstrated envelope, and several (glass doors,
        corridors with no indicator objects) are known hard cases for the
        current perception vocabulary.
  \item \textbf{Validation data overlaps the training distribution.} As noted in
        Section~\ref{sec:door_eval}, the door detector's held-out split shares
        its source datasets with training; cross-domain performance is
        characterized only by the live trials.
\end{enumerate}